\documentclass[10pt]{article}
\usepackage[a4paper,margin=22mm]{geometry}
\usepackage{amsmath,amssymb,booktabs,xcolor}
\setlength{\parskip}{0.35em}\setlength{\parindent}{0pt}
\newcommand{\Wz}{W\!z_{CoM}}
\title{\vspace{-16mm}\textbf{Drop-in text for Sec.~VI-E (deviation bound)
and Sec.~VI-B (ramp tracking)}\vspace{-3mm}}
\date{\small numbering continues from Eq.~(86); symbols follow the manuscript}
\begin{document}\maketitle

\section*{VI-E continuation: impulse response and the bound on $e_\varphi$}

\emph{After Proposition 1 / Eq.~(86).}

\medskip
The homogeneous part of (86) is spanned by $\cosh C_2\tau$ and
$\sinh C_2\tau$ with $C_2 = \sqrt{\Wz/J_P}$ from (20).  Its impulse
response --- the solution of $J_P\ddot h - \Wz\,h = \delta(\tau)$ from
rest --- follows from the momentum jump
$J_P[\dot h(0^+)-\dot h(0^-)] = 1$, the position being unable to jump:
\begin{equation}
  h_\varphi(\tau) = \frac{\sinh(C_2\tau)}{J_P C_2},
  \qquad h_\varphi(0) = 0,
  \qquad \dot h_\varphi(0) = \frac{1}{J_P}.
  \tag{87}
\end{equation}
Because $h_\varphi(0) = 0$, the boundary term generated by
differentiating a convolution with a variable upper limit,
\[
  \frac{d}{d\tau}\int_0^\tau h_\varphi(\tau-s)\rho(s)\,ds
  = \underbrace{h_\varphi(0)}_{=\,0}\rho(\tau)
  + \int_0^\tau \dot h_\varphi(\tau-s)\rho(s)\,ds ,
\]
vanishes for \emph{any} forcing.  The zero initial conditions in (86)
therefore set both homogeneous coefficients to zero --- no assumption on
$\rho(0)$ is required --- and the deviation is the Duhamel integral
alone,
\begin{equation}
  e_\varphi(\tau) = \int_0^\tau h_\varphi(\tau-s)\,\rho(s)\,ds ,
  \qquad
  e_\omega(\tau) = \int_0^\tau \dot h_\varphi(\tau-s)\,\rho(s)\,ds .
  \tag{88}
\end{equation}

\paragraph{The two channels.}
Two terms of (79) leave the estimator span: the gravity remainder (83)
and the bilinear ground-effect coupling (79)(iv),
\begin{equation}
  \rho(\tau) = \rho_\varphi(\tau) + \rho_{GE}(\tau),
  \qquad
  \rho_{GE}(\tau) = \beta_M \dot M \tau\varphi .
  \tag{89}
\end{equation}
Both vanish at the onset --- with $\varphi_{\rm nom}\sim\tau^3$ the
gravity remainder is $\mathcal{O}(\varphi^2)\sim\tau^6$ and the
bilinear term is $\tau\cdot\tau^3 = \tau^4$ --- so each is largest at
the window end.

\paragraph{An envelope bound that can be verified by hand.}
Replacing $\rho$ by a constant $\bar\rho$ and integrating the two
kernels, $\int_0^\tau h_\varphi = (\cosh C_2\tau - 1)/(J_PC_2^2)$ and
$\int_0^\tau \dot h_\varphi = \sinh(C_2\tau)/(J_PC_2)$, and using
$J_PC_2^2 = \Wz$,
\begin{equation}
  \bigl|e_\varphi(\tau_{\rm end})\bigr| \le
  \frac{\bar\rho}{\Wz}\bigl(\cosh x - 1\bigr),
  \qquad
  \bigl|e_\omega(\tau_{\rm end})\bigr| \le
  \frac{\bar\rho\,C_2}{\Wz}\,\sinh x ,
  \qquad x \triangleq C_2\tau_{\rm end}.
  \tag{90}
\end{equation}
The inertia cancels in both: $J_PC_2 = \Wz/C_2$ likewise. Only $\Wz$,
$C_2$, the dimensionless group $x$ and the forcing level enter.

\paragraph{The window length is bounded a priori by the tilt cap.}
Neither channel maximum nor the kernel factor requires $x$ to be
solved for. With the nominal excursion of (20),
\begin{equation}
  \varphi_{\rm end} = \frac{\dot M}{\Wz C_2}\,\Lambda(x),
  \qquad \Lambda(x) \triangleq \sinh x - x
  = \sum_{k\ge1}\frac{x^{2k+1}}{(2k+1)!} .
  \tag{91}
\end{equation}
Every term of $\Lambda$ is positive, so $\Lambda(x)\ge x^3/6$, and
since $\Lambda$ increases, imposing the tilt cap
$\varphi_{\rm end}\le\varphi_{\max}$ inverts in closed form:
\begin{equation}
  x \le \bar x \triangleq
  \Bigl(\frac{6\,\varphi_{\max}\Wz C_2}{\dot M}\Bigr)^{1/3},
  \qquad
  \tau_{\rm end} \le \frac{\bar x}{C_2}
  = \Bigl(\frac{6\,\varphi_{\max}\Wz}{C_2^{2}\dot M}\Bigr)^{1/3},
  \tag{92}
\end{equation}
where $J_P = \Wz/C_2^{2}$ has been used to remove the inertia. The
retained term is the inertia-only rise
$\varphi = \dot M\tau^3/(6J_P)$ obtained by dropping the destabilising
$\Wz\varphi$: it is the slowest admissible ascent and therefore reaches
the cap latest, which is why (92) bounds $x$ from above and why the
worst case is the \emph{slowest} ramp, not the fastest. The moment
increment over the window and the coupling channel follow with no
measured quantity at all,
\begin{equation}
  \Delta M_{\rm win} = \dot M\tau_{\rm end}
  \le \Bigl(\frac{6\,\varphi_{\max}\Wz\dot M^{2}}{C_2^{2}}\Bigr)^{1/3},
  \qquad
  \rho_{GE,\max} = |\beta_M|\,\Delta M_{\rm win}\varphi_{\rm end}
  \le |\beta_M|\,\varphi_{\max}
  \Bigl(\frac{6\,\varphi_{\max}\Wz\dot M^{2}}{C_2^{2}}\Bigr)^{1/3},
  \tag{93}
\end{equation}
and the kernel factors of (90) close on the same constant through the
identity $\cosh x - 1 = \Lambda(x) + (x - 1 + e^{-x})$ with
$0\le x-1+e^{-x} < x$, together with $\sinh x = \Lambda(x) + x$:
\begin{equation}
  \cosh x - 1 \;\le\; \sinh x \;\le\; \Lambda(x) + \bar x
  \;\le\; \frac{\bar x^{3}}{6} + \bar x .
  \tag{94}
\end{equation}
Both kernels therefore share one bracket.\footnote{%
The step matters. $\Lambda(x)$ is known \emph{as an equality} --- the
tilt cap supplies it --- so only the residual linear term needs $\bar
x$. Substituting $\bar x$ into $\sinh$ instead is exponentially lossy:
at the slowest ramp $\bar x = 7.23$ gives $\sinh\bar x = 691$ against
$\Lambda + \bar x = 70.3$, a factor of nine thrown away for nothing.}

\paragraph{Why the envelope alone does not settle the question.}
Nothing beyond $|\rho|\le\rho_{\max}$ enters (90): no monotonicity, no
shape, no property of the trajectory but the window length. That
generality is also its cost. A constant envelope spreads the forcing
over the whole window, whereas both channels are concentrated at its
end ($\tau^6$ and $\tau^4$), and the mismatch grows with $x$, i.e.\
with the \emph{slowest} ramps. Evaluated with $\bar\rho = \rho_{\max}$
over the identified constants, (90) returns a relative deviation of
$467\%$ at $\dot M = 0.10$~N\,m/s --- larger than the excursion it
bounds. The envelope form is a sanity check on the algebra, not a
result.

What restores it is the one structural fact available for free: $\rho$
rises monotonically from zero at the onset while the kernel
$\cosh\bigl(C_2(\tau-s)\bigr)$ decreases in $s$. For functions of
\emph{opposite} monotonicity Chebyshev's integral inequality runs in
the direction needed, and each maximum in (90) may be replaced by the
channel's time average. With the nominal shapes
$\rho_\varphi\propto\Lambda(u)^2$ and $\rho_{GE}\propto u\,\Lambda(u)$,
$u = C_2\tau$,
\begin{equation}
  \bar\rho = R_\varphi(x)\,\rho_{\varphi,\max}
           + R_{GE}(x)\,\rho_{GE,\max},
  \qquad
  R_\varphi \le \tfrac17, \qquad R_{GE} \le \tfrac15,
  \tag{95}
\end{equation}
the limits being the $x\to0$ values where the two channels are sixth
and fourth order; over the flown $x\in[1.79,5.20]$ they are
$0.089$--$0.133$ and $0.145$--$0.191$. The same evaluation that gave
$467\%$ returns $43.7\%$ with (95) in place.

\paragraph{Relative form: the result, free of $x$.}
Dividing (90) by the nominal excursion
$\varphi_{\rm nom} = (C_1/C_2)\Lambda(x)$ with $C_1 = \dot M/\Wz$ and
writing $\Delta M_{\rm win} = \dot M\tau_{\rm end} = \dot M x/C_2$,
\begin{equation}
  \frac{|e_\varphi|_{\rm end}}{\varphi_{\rm nom,end}}
  \le \frac{\Psi(x)\,\bar\rho}{\Delta M_{\rm win}},
  \qquad
  \Psi(x) = \frac{x(\cosh x - 1)}{\sinh x - x} .
  \tag{96}
\end{equation}
Here the two approximations work against each other: $\Psi$ grows with
$x$ and the averages of (95) shrink with it, and each product is
bounded uniformly --- $\Psi R_\varphi$ rises monotonically from $3/7$
to $1/2$ and $\Psi R_{GE}$ from $3/5$ to $1$. Hence
\begin{equation}
  \boxed{\;\frac{|e_\varphi|_{\rm end}}{\varphi_{\rm nom,end}}
  \;\le\; \frac{\tfrac12\rho_{\varphi,\max} + \rho_{GE,\max}}
               {\Delta M_{\rm win}},
  \qquad
  |\Delta M_{\rm crit}| \;\le\; \frac{\rho_{\varphi,\max}}{7}
                              + \frac{\rho_{GE,\max}}{5}\;}
  \tag{97}
\end{equation}
independently of $x$, $C_2$, $J_P$ and $\Wz$: the window length that
(92) went to some trouble to bound does not survive into the
conclusion, and is needed only to bound $\rho_{GE,\max}$ through (93).
The second inequality follows because the onset minimises the
two-segment residual: linearising its stationarity condition gives
$\Delta M_{\rm crit} = -\!\int_0^{\tau_{\rm end}}\!\rho\,w$ with a
non-negative weight normalised to $\int w = 1$, so
$|\Delta M_{\rm crit}|\le\bar\rho$. A reviewer can evaluate (97) from
numbers already in the paper.

\paragraph{No pivot inertia is assumed anywhere.}
$C_2$ is not computed from $J_P$; the two-stage calibration of Sec.~V
identifies the pair $(C_2, K)$ directly from the onset fits, and
$\Wz = 1/K$ is one of them. $J_P = 1/(KC_2^{2})$ is a derived
by-product that enters none of (90)--(97) --- which is as well, since
the identified $J_P$ spans $0.10$--$0.28$~kg\,m$^2$ against the rigid
parallel-axis value $J_{CoM} + md^2 \approx 0.43$~kg\,m$^2$, a gap the
compliant contact makes uninterpretable as an inertia. Over the ten
configurations the identified constants span
$C_2\in[3.52, 5.67]$~s$^{-1}$ and $\Wz\in[2.96, 6.65]$~N\,m, so with
$\varphi_{\max} = 10^\circ$ and $\dot M\in[0.10,1.20]$~N\,m/s, (92)
gives $\bar x\le 7.23$ at the slowest ramp and (93) gives
$\rho_{GE,\max}\le 4.5$~mN\,m at the fastest --- the two extremes lie
at opposite corners of the box --- against
$\rho_{\varphi,\max}\le\tfrac12 W(l_p+\lambda_{\rm off})
\varphi_{\max}^{2} = 77.0$~mN\,m (roll; $62.5$ pitch), or $50.7$~mN\,m
if the stabilising $-\Wz\sin\varphi$ term of $G''$ is retained at the
cap. Because $\bar x$ depends on the constants only as
$(\Wz C_2/\dot M)^{1/3}$, a factor-of-two error in $\Wz$ moves it by
$26\%$: the chain is least sensitive to exactly the constants that are
least certain.

\paragraph{Evaluation.}
Every quantity must be read off the \emph{same} window:
$\rho_{\varphi,\max}$ and $\rho_{GE,\max}$ at the excursion that window
reaches, $\Delta M_{\rm win}$ as its own moment increment.  Over the
$140$ gate-passing runs, with $W = 31.59$~N,
$l_p+\lambda_{\rm off} = 0.160/0.130$~m (roll/pitch), $z_{CoM} = 0.30$~m
and $\beta_M = 0.0345$~rad$^{-1}$:

\begin{center}\small
\setlength{\tabcolsep}{4pt}
\begin{tabular}{@{}lrrrrrr@{}}
\toprule
$\dot M$ [N\,m/s] & $0.10$ & $0.30$ & $0.65$ & $1.20$ & median & worst \\
\midrule
$|\varphi|_{\rm end}$ [deg]          & $5.23$ & $5.90$ & $6.55$ & $6.17$ & $5.95$ & $9.33$\\
$\Delta M_{\rm win}$ [N\,m]          & $0.072$& $0.165$& $0.306$& $0.479$& $0.221$& ---\\
$\rho_{\varphi,\max}$ [mN\,m]        & $19.7$ & $22.5$ & $30.4$ & $27.2$ & $23.6$ & $58.3$\\
$\rho_{GE,\max}$ [mN\,m]             & $0.24$ & $0.59$ & $1.24$ & $1.85$ & $0.68$ & $3.00$\\
$x = C_2\tau_{\rm end}$              & $3.47$ & $2.64$ & $2.15$ & $1.79$ & $2.67$ & $5.20$\\
\midrule
$|e_\varphi|/\varphi_{\rm nom}$ [\%] & $13.8$ & $6.0$  & $4.7$  & $2.9$  & $4.9$  & $25.1$\\
$|e_\varphi|$ [deg]                  & $0.76$ & $0.34$ & $0.31$ & $0.18$ & $0.32$ & $1.68$\\
$|\Delta M_{\rm crit}|$ [mN\,m]      & $1.94$ & $2.73$ & $3.77$ & $3.76$ & $3.01$ & $7.86$\\
$|\Delta\lambda_{\rm off}|$ [mm]     & $0.061$& $0.086$& $0.119$& $0.119$& $0.095$& $0.249$\\
\midrule
ground-effect share of $\bar\rho$ [\%] & $1.8$ & $4.0$ & $6.0$ & $8.8$ & $5.1$ & $25.6$\\
\bottomrule
\end{tabular}
\end{center}

The table entry is the per-run propagation, which convolves the
measured $\rho$ with the exact kernel and needs no envelope at all.
Using instead the $x$-free constants of (97) the same runs give
$|\Delta M_{\rm crit}|\le 8.8$~mN\,m, i.e.\
$|\Delta\lambda_{\rm off}|\le 0.28$~mm --- the price of not evaluating
$R_\varphi, R_{GE}$ is under $12\%$.  Applied a priori at the tilt cap
rather than at the flown excursions, (97) gives $11.9$~mN\,m and
$0.38$~mm, which is the guarantee the method carries into a
configuration it has not seen.

The bilinear ground-effect channel contributes $5.1\%$ of $\bar\rho$ in
the median and $25.6\%$ at worst; it grows with the ramp rate because
$\rho_{GE}\propto\dot M\tau$, whereas the gravity remainder depends on
the excursion alone.  \textcolor{blue}{The threshold bound
$|\Delta\lambda_{\rm off}| \le 0.25$~mm is the figure that matters: it
does not involve $\Delta M_{\rm win}$, and so does not inherit the loss
the relative normalisation suffers at slow ramps.}  The relative bound
is loosest at $0.10$~N\,m/s, where the window closes on a small moment
increment while the excursion is unchanged --- the same ordering the
excitation design reports --- and it is conservative by roughly a factor
$5$ overall, since it forwards the \emph{unprojected} remainder
($23.6$--$58.3$~mN\,m here, against $1.2$--$12.3$~mN\,m measured after
the estimator absorbs its own degrees of freedom).

\section*{VI-B: ramp input tracking (rate limits)}

The rotor response identified in Sec.~V (Table~4:
$\omega_{n,\rm rot} = 21.74$~rad/s, $\zeta = 0.7814$,
$\tau_d = 15.67$~ms) tracks a ramp with the steady-state lag
\[
  t_{\rm lag} = \frac{2\zeta}{\omega_{n,\rm rot}} + \tau_d
  = 71.9 + 15.7 = 87.6~\text{ms},
\]
so the realised moment trails the command by
$\dot M\,t_{\rm lag}$, i.e.\ $8.8$~mN\,m at $0.10$~N\,m/s and
$105$~mN\,m at $1.20$~N\,m/s.

This lag does \emph{not} bias $M_{\rm crit}$.  The identification
consumes the \emph{realised} moment, reconstructed from the measured
rotor speeds, not the commanded one; a common lag therefore shifts the
moment trace and the rate trace together and leaves the moment read at
the detected onset unchanged.  What the lag does consume is window
length, which the sample-count gate enforces.

Two residual effects remain and are both gated.  First, the realised
\emph{slope} may differ from the commanded rate; the run-level gate
rejects $|\dot M_{\rm meas} - \dot M_{\rm cmd}|/\dot M_{\rm cmd} > 3\%$,
and since $C_1 = K\dot M$, a $3\%$ rate error perturbs the amplitude
constant by $3\%$ --- well inside the $\pm20\%$ mis-specification the
sensitivity analysis shows to be harmless.  Second, the start of the
ramp is not linear while the rotor transient decays; with
$4/(\zeta\omega_{n,\rm rot}) = 235$~ms to $2\%$, that transient is
confined to the low-moment part of the window, which the slope gate
excludes by scoring only the segment above the per-axis moment floor.
The linearity gate ($30$~mN\,m RMSE about the best-fit line) then acts
as a fault detector for ramps that are stepped or aborted outright.

\end{document}
